
\SGAProof{Proof of \SGARefLink{sga3:xv:theorem:8-8:item:a}{Theorem 8.8(a)}}
If the functor \(\mathcal T_C\) is representable, it commutes with
adic limits of artinian rings. Consequently, by \SGARefLink{sga3:xv:remark:8:2}{8.2} and
Proposition 8.5(b), \(G\) has property \(\mathrm{ATC}\). To prove the
converse, we shall use the following result, which will be proved in
EGA VI and may also be found in Murre's exposé,
\emph{Séminaire Bourbaki}, May 1965, no.~294, Theorem~1,
Corollary~2.

\SGAStatement{\SGAInlineRefTarget{sga3:xv:lemma:8-10}Lemma 8.10}{Let \(S\) be a locally noetherian
prescheme and let \(F\) be a contravariant functor from
\(\operatorname{Sch}/S\) to the category of sets. In order that
\(F\) be representable by an \(S\)-prescheme étale and separated
over \(S\), it is necessary and sufficient that \(F\) satisfy the
following five properties:
\begin{enumerate}
\renewcommand{\labelenumi}{\textup{\roman{enumi})}}
\item \(F\) is a sheaf for the fpqc topology (Exposé IV).
\item \(F\) commutes with filtered inductive limits of rings
(Exposé \SGARefLink{sga3:xi:remark:3:2}{XI~3.2}).
\item \(F\) commutes with adic limits of local artinian rings
(\SGARefLink{sga3:xv:proposition:8-1}{8.1}\textup{(i)}).
\item \(F\) satisfies the valuative criterion for separatedness:
for every \(S\)-prescheme \(S'\) that is the spectrum of a discrete
valuation ring, with generic point \(t\), the canonical map
\[
  F(S')\longrightarrow F(t)
\]
is injective.
\item \(F\) is infinitesimally étale (Exposé \SGARefLink{sga3:xi:remark:1:8}{XI~1.8}).
\end{enumerate}}
\SGARefTarget{sga3:xv:lemma:8:10}

We show that the functor \(\mathcal T_C\) of \SGARefLink{sga3:xv:theorem:8-8}{Theorem 8.8} satisfies
the five conditions of \SGARefLink{sga3:xv:lemma:8-10}{Lemma 8.10}.

\begin{enumerate}
\renewcommand{\labelenumi}{\textup{\roman{enumi})}}
\item The functor \(\mathcal T\), and likewise \(\mathcal T_C\), is
a sheaf for the fpqc topology whenever \(G\) is of finite
presentation over \(S\). Indeed, every monomorphism from a torus
into \(G\) is then an immersion (Exposé \SGARefLink{sga3:viii:proposition:7:9}{VIII~7.9}), and the assertion
follows from fpqc descent for subpreschemes.

\item \SGARefLink{sga3:xv:corollary:6-3}{Corollary 6.3} proves that \(\mathcal T\) commutes with filtered
inductive limits of rings when \(G\) is of finite presentation over
\(S\); the same property for \(\mathcal T_C\) follows immediately.

\item By \SGARefLink{sga3:xv:proposition:8:5}{Proposition 8.5}, condition \textup{(iii)} is precisely
equivalent to property \(\mathrm{ATC}\).

\item If \(S\) is the spectrum of a discrete valuation ring, two
sub-tori of \(G\) with the same generic fibre coincide. More
precisely, each is the identity component of the schematic closure
in \(G\) of that common generic fibre.

\item This follows from \SGARefLink{sga3:xv:corollary:2:3}{Proposition 2.3} because \(G\) is flat over
\(S\).
\end{enumerate}

\SGAProof{Proof of \SGARefLink{sga3:xv:theorem:8-8:item:b}{Theorem 8.8(b)}}
Proceeding as in Exposé \SGARefLink{sga3:xi:corollary:4:2}{XI~4.2}, it is enough to prove that the
product group \(T\times_S G\) again has property
\(\mathrm{ATC}\), which is immediate from the definition.

\SGAProof{Proof of \SGARefLink{sga3:xv:theorem:8-9:item:a}{Theorem 8.9(a)}}
After replacing \(G\), if necessary, by its identity component
(Exposé VI\textsubscript{B}~\SGARefLink{sga3:xv:lemma:3:10}{3.10}), we may suppose that \(G\) has
connected fibres. If \(T\) is a sub-torus of \(G\), its centralizer
\(\operatorname{Centr}_G(T)\) is represented (Exposé \SGARefLink{sga3:xi:corollary:6:11}{XI~6.11}) by a
subgroup prescheme of \(G\), smooth over \(S\)
(Exposé \SGARefLink{sga3:xi:corollary:2:4}{XI~2.4}), with connected fibres
(Exposé \SGARefLink{sga3:xii:theorem:6-6:item:b}{XII~6.6(b)}). It is therefore an element of
\(\mathcal{CT}(S)\), where \(\mathcal{CT}\) is the functor defined
in \SGARefLink{sga3:xv:proposition:7-1:item:i}{7.1(i)}. The rule
\[
  T\longmapsto\operatorname{Centr}_G(T)
\]
clearly defines an \(S\)-morphism
\[
  u:\mathcal T\longrightarrow\mathcal{CT}.
\]
Since \(\mathcal{CT}\) is represented by an \(S\)-prescheme smooth
and quasi-projective over \(S\) (\SGARefLink{sga3:xv:proposition:7-1:item:i}{7.1(i)}), it remains only to prove
that \(u\) is representable by a separated étale morphism.

After a suitable base change \(S'\to S\), with
\(S'=\mathcal{CT}\) and hence \(S'\) locally noetherian, we are
reduced to the following problem. Let \(S\) be a locally noetherian
prescheme, let \(G\) be an \(S\)-group prescheme smooth and of
finite type over \(S\), with connected fibres, and let \(H\) be a
subgroup of \(G\), smooth over \(S\), with connected fibres.
Consider the subfunctor \(X\) of \(\mathcal T_{C,H}\) such that,
for every \(S\)-prescheme \(S'\), \(X(S')\) is the set of central
sub-tori \(T\) of \(H_{S'}\) for which
\[
  \operatorname{Centr}_{G_{S'}}(T)=H_{S'}.
\]
We shall show that \(X\) is represented by an \(S\)-prescheme
étale and separated over \(S\).

Indeed, \(G\) has property \(\mathrm{AT}\) by hypothesis, and so
does \(H\) by Proposition 8.7(ii). Since \(H\) has connected
fibres, it also has property \(\mathrm{ATC}\) by
Proposition 8.7(i). Moreover, \(H\) is smooth over \(S\), hence
flat. Theorem 8.8(a) therefore represents \(\mathcal T_{C,H}\) by
an \(S\)-prescheme étale and separated over \(S\). It is enough to
show that the canonical monomorphism
\[
  X\longrightarrow\mathcal T_{C,H}
\]
is represented by an open immersion.

Put \(S'=\mathcal T_{C,H}\), and let \(K\) be the centralizer in
\(G_{S'}\) of the universal central torus of \(H_{S'}\). The group
\(K\) is a smooth group scheme over \(S'\), with connected fibres,
and contains \(H_{S'}\). By definition,
\[
  X=\prod_{K/S'} H_{S'}/K,
\]
which is represented by the open-and-closed subprescheme of \(S'\)
over which \(H_{S'}\) and \(K\) have the same relative dimension.

The proof of \SGARefLink{sga3:xv:theorem:8-9:item:b}{Theorem 8.9(b)} is analogous to that of
\SGARefLink{sga3:xv:theorem:8-8:item:b}{Theorem 8.8(b)}.

\SGAStatement{Corollary 8.11}{Let \(S\) be a prescheme and let
\(G\) be an \(S\)-group prescheme smooth and of finite presentation
over \(S\). If the abelian rank of the fibres of \(G\) is a locally
constant function on \(S\)---in particular, if the fibres of \(G\)
are affine---then the functor of sub-tori of \(G\) is represented
by an \(S\)-prescheme smooth and separated over \(S\). The same
holds for the functor
\(\operatorname{Hom}_{S\text{-gr}}(T,G)\), for every \(S\)-torus
\(T\).}
\SGARefTarget{sga3:xv:corollary:8:11}

The assertion is local on \(S\), so we may suppose that \(S\) is
affine and that the abelian rank of the fibres of \(G\) is constant.
There are then a noetherian affine scheme \(S_0\) and an
\(S_0\)-group prescheme \(G_0\), smooth and of finite type over
\(S_0\), such that
\[
  G\simeq G_0\times_{S_0}S
\]
as \(S\)-group preschemes (Exposé VI\textsubscript{B}, \S10).
The abelian rank of the fibres of an \(S\)-group prescheme of finite
presentation is an ind-constructible function (6.3 bis). A standard
argument (EGA IV~8) therefore allows us, in the present situation,
to choose \(G_0\) so that the abelian rank of its fibres is constant
on \(S_0\). Proposition 8.7(iii) gives property \(\mathrm{AT}\) for
\(G_0\), and \SGARefLink{sga3:xv:theorem:8:8}{Theorem 8.9} then represents its functor of sub-tori by
an \(S_0\)-prescheme smooth and separated over \(S_0\). Base change
gives the assertion for \(G\). The argument for
\(\operatorname{Hom}_{S\text{-gr}}(T,G)\) is analogous.

\medskip\noindent\textbf{Generalization of \SGARefLink{sga3:xv:theorem:8:8}{Theorem 8.9}.}
The functor \(\mathcal T\) of sub-tori of a smooth group \(G\) need
not be representable. We now give sufficient conditions for a
subfunctor of \(\mathcal T\) to be representable.

\SGAStatement{\SGAInlineRefTarget{sga3:xv:proposition:8-12}Proposition 8.12}{Let \(S\) be a locally noetherian
prescheme, let \(G\) be an \(S\)-group prescheme smooth over \(S\)
with connected fibres, and let \(F\) be an \(S\)-subfunctor of the
functor \(\mathcal T\) of sub-tori of \(G\). Suppose that:
\begin{enumerate}
\renewcommand{\labelenumi}{\textup{\roman{enumi})}}
\item\SGAInlineRefTarget{sga3:xv:proposition:8-12:item:i} \(F\) is a sheaf for the fpqc topology (Exposé IV);
\item \(F\) commutes with filtered inductive limits of rings
(Exposé \SGARefLink{sga3:xi:remark:3:2}{XI~3.2});
\item \(F\) commutes with adic limits of local artinian rings
(Exposé \SGARefLink{sga3:xi:remark:3:3}{XI~3.3});
\item\SGAInlineRefTarget{sga3:xv:proposition:8-12:item:iv} \(F\) is infinitesimally smooth over \(S\)
(Exposé \SGARefLink{sga3:xi:remark:1:8}{XI~1.8});
\item\SGAInlineRefTarget{sga3:xv:proposition:8-12:item:v} \(F\) is stable under inner automorphisms of \(G\): for every
\(S\)-prescheme \(S'\),
\[
  T\in F(S'),\quad g\in G(S')
  \quad\Longrightarrow\quad
  \operatorname{int}(g)T\in F(S').
\]
\end{enumerate}
Then \(F\) is represented by an \(S\)-prescheme smooth and
separated over \(S\).}
\SGARefTarget{sga3:xv:proposition:8:12}

\SGAStatement{Proposition 8.12 bis}{Let \(S\) and \(G\) be as in
\SGARefLink{sga3:xv:proposition:8-12}{Proposition 8.12}, let \(T\) be an \(S\)-torus, and let \(F\) be a
subfunctor of \(\operatorname{Hom}_{S\text{-gr}}(T,G)\). Suppose
that \(F\) satisfies conditions \textup{(i)}--\textup{(iv)} of
\SGARefLink{sga3:xv:proposition:8-12}{Proposition 8.12} and is stable under inner automorphisms of \(G\):
for every \(S\)-prescheme \(S'\),
\[
  u\in F(S'),\quad g\in G(S')
  \quad\Longrightarrow\quad
  \operatorname{int}(g)u\in F(S').
\]
Then \(F\) is represented by an \(S\)-prescheme smooth and
separated over \(S\).}
\SGARefTarget{sga3:xv:proposition:8:12:bis}

\SGAProof{Proof of \SGARefLink{sga3:xv:proposition:8-12}{Proposition 8.12}}
The proof of \SGARefLink{sga3:xv:proposition:8:12}{Proposition 8.12 bis} is entirely analogous and is left
to the reader. Let
\[
  u:F\longrightarrow\mathcal{CT}
\]
be the \(S\)-morphism which assigns to every
\(T\in F(S')\) the element
\(\operatorname{Centr}_{G_{S'}}(T)\) of
\(\mathcal{CT}(S')\) (\SGARefLink{sga3:xv:proposition:7-1:item:i}{7.1(i)}). Since \(\mathcal{CT}\) is represented
by an \(S\)-prescheme smooth and quasi-projective over \(S\), we
are reduced to proving that \(u\) is representable.

After the base change \(\mathcal{CT}\to S\), the problem is the
following. Given \(S\) and \(G\) as above and a smooth subgroup
\(H\) of \(G\) with connected fibres, represent the functor \(F_H\)
defined by
\[
  F_H(S')=
  \left\{T\in F(S')\ \middle|\
  H_{S'}=\operatorname{Centr}_{G_{S'}}(T)\right\}.
\]
We show that \(F_H\) is étale and separated over \(S\) by verifying
the five conditions of \SGARefLink{sga3:xv:lemma:8-10}{Lemma 8.10}.

The first three conditions follow directly from
\SGARefLink{sga3:xv:proposition:8-12:item:i}{Proposition 8.12(i)}--(iii). We have already observed that
\(\mathcal T_{C,H}\) is separated and infinitesimally étale.
Thus \(F_H\), being a subfunctor of \(\mathcal T_{C,H}\), is
separated and infinitesimally unramified (Exposé \SGARefLink{sga3:xi:remark:1:8}{XI~1.8}). It remains
to prove that \(F_H\) is infinitesimally smooth.

Let \(S\) be the spectrum of a local artinian ring, let
\(S_0\subseteq S\) be defined by a nilpotent ideal, and let
\(T_0\in F_H(S_0)\). By \SGARefLink{sga3:xv:proposition:8-12:item:iv}{Proposition 8.12(iv)}, \(T_0\) lifts to an
element \(T'\in F(S)\). Since \(H\) is smooth, \(T_0\) also lifts
to a sub-torus \(T''\) of \(H_S\) (Exposé \SGARefLink{sga3:ix:theorem:3:6:bis}{IX~3.6 bis}). This
sub-torus is conjugate to \(T'\) by an element \(g\in G(S)\)
(\emph{loc. cit.}), and hence \(T''\in F(S)\) by
\SGARefLink{sga3:xv:proposition:8-12:item:v}{Proposition 8.12(v)}. The centralizer
\(\operatorname{Centr}_{G_S}(T'')\) is smooth over \(S\) and agrees
with \(H_{S_0}\) over \(S_0\). Therefore \(T''\) lies in the centre
of \(H_S\) (Exposé \SGARefLink{sga3:ix:theorem:5:6:item:a}{IX~5.6(a)}), and its centralizer in \(G\) equals
\(H_S\). Thus \(T''\in F_H(S)\), and we take \(T=T''\).

\SGAStatement{Corollary 8.13}{Let \(S\) be a prescheme and let
\(G\) be an \(S\)-group prescheme smooth and of finite presentation
over \(S\), with connected fibres. Let \(\mathcal T_D\) be the
subfunctor of \(\mathcal T\) for which \(\mathcal T_D(S')\) is the
set of sub-tori \(T\) of \(G_{S'}\) such that, for every
\(s'\in S'\), \(T_{s'}\) is contained in the derived subgroup of
\(G_{s'}\) (Exposé VI\textsubscript{B}~7). Then \(\mathcal T_D\)
is represented by an \(S\)-prescheme smooth and separated over
\(S\).}
\SGARefTarget{sga3:xv:corollary:8:13}

\SGAStatement{Corollary 8.13 bis}{Let \(S\) and \(G\) be as in
\SGARefLink{sga3:xv:corollary:8:13}{Corollary 8.13}, and let \(T\) be an \(S\)-torus. Let
\[
  \operatorname{Hom}^{\mathrm{der}}_{S\text{-gr}}(T,G)
\]
be the subfunctor of \(\operatorname{Hom}_{S\text{-gr}}(T,G)\)
whose \(S'\)-points are the \(S'\)-morphisms
\(u:T_{S'}\to G_{S'}\) such that, for every \(s'\in S'\), the
morphism \(u_{s'}\) factors through the derived group of
\(G_{s'}\). Then this functor is represented by an \(S\)-prescheme
smooth and separated over \(S\).}
\SGARefTarget{sga3:xv:corollary:8:13:bis}

The two corollaries have analogous proofs; we prove
\SGARefLink{sga3:xv:corollary:8:13:bis}{Corollary 8.13 bis}. By the usual procedure, we reduce to the case
in which \(S\) is noetherian. We verify the five conditions of
\SGARefLink{sga3:xv:proposition:8:12}{Proposition 8.12 bis}.

Conditions \textup{(i)} and \textup{(iv)} follow immediately from
the corresponding properties of the functor
\[
  \operatorname{Hom}_{S\text{-gr}}(T,G).
\]
Condition \textup{(v)} holds because the derived group of an
algebraic group is invariant (Exposé VI\textsubscript{B}, \S7).

For condition \textup{(ii)}, a standard reduction (EGA IV~8)
reduces us to the following assertion. Let \(S\) be an integral
noetherian scheme with generic point \(\eta\), and let
\(u:T\to G\) be an \(S\)-group morphism whose generic fibre factors
through the derived group of \(G_\eta\). Then there is a
neighbourhood \(U\) of \(\eta\) such that, for every \(s\in U\),
\(u_s\) factors through the derived group of \(G_s\). This is an
immediate consequence of Exposé VI\textsubscript{B}~\SGARefLink{prop:sga3-VIB-10.12}{10.12}.

For condition \textup{(iii)}, retain the notation of \SGARefLink{sga3:xv:proposition:4:3}{Proposition
4.3}. We must show that an element
\[
  (u_m)_{m\in\mathbb N}\in
  \varprojlim_m
  \operatorname{Hom}^{\mathrm{der}}_{S_m\text{-gr}}(T_m,G_m)
\]
is admissible in the sense of \SGARefLink{sga3:xv:proposition:4:3}{Proposition 4.3} and comes from an
element of
\(\operatorname{Hom}^{\mathrm{der}}_{S\text{-gr}}(T,G)\).
As in \SGARefLink{sga3:xv:proposition:8-1}{8.1}, we reduce immediately to the case in which \(S\) is the
spectrum of a complete discrete valuation ring. Let \(t\) be its
generic point, \(s\) its closed point, and let \(D\) be the
schematic closure in \(G\) of the derived group \(D_t\) of \(G_t\).

\begin{enumerate}
\renewcommand{\labelenumi}{\textup{\alph{enumi})}}
\item The group \(D_s\) contains the derived group of \(G_s\).
Indeed, \(D_t\) is invariant in \(G_t\), and \(G\) is flat over
\(S\), so \(D\) is invariant in \(G\). Moreover, the morphism
\[
  G\times_S G\longrightarrow G,
  \qquad (x,y)\longmapsto xyx^{-1}y^{-1},
\]
factors through \(D_t\) on the generic fibre and therefore factors
through \(D\). Consequently \(G_s/D_s\) is commutative.

\item If
\(u_m\in
\operatorname{Hom}^{\mathrm{der}}_{S_m\text{-gr}}(T_m,G_m)\),
then \(u_m\) factors through \(D_m\). Its closed fibre \(u_s\)
factors through the derived group of \(G_s\), hence through \(D_s\)
by \textup{(a)}. Since \(D_m\) is flat over \(S_m\) and invariant
in \(G_m\), the quotient
\[
  H_m=G_m/D_m
\]
is representable (Exposé VI\textsubscript{A}, \S4). The image of
\(T_m\) in \(H_m\) is a torus (Exposé \SGARefLink{sga3:ix:theorem:6:8}{IX~6.8}) with trivial closed
fibre, hence is trivial. Thus \(u_m\) factors through \(D_m\).

\item The family \((u_m)_{m\in\mathbb N}\) is admissible and lifts
to a morphism \(T\to G\) belonging to
\(\operatorname{Hom}^{\mathrm{der}}_{S\text{-gr}}(T,G)\).
By what precedes and \SGARefLink{sga3:xv:theorem:4-1-bis}{Proposition 4.1 bis}, it is enough to prove that
\(D_t\) is affine. This follows from the next lemma.
\end{enumerate}

\SGAStatement{Lemma 8.14}{Let \(G\) be a smooth connected
algebraic group over a field \(k\). Then the derived group \(D\) of
\(G\) is affine.}
\SGARefTarget{sga3:xv:lemma:8:14}

Formation of \(D\) commutes with extension of the base field
(Exposé VI\textsubscript{B}~7), so we may suppose that \(k\) is
algebraically closed. The group \(G\) is then an extension of an
abelian variety \(A\) by a linear group \(L\). Since \(A\) is
commutative, \(D\) is necessarily contained in \(L\), and is
therefore affine.

\medskip\noindent\textbf{Maximal tori.}

\SGAStatement{\SGAInlineRefTarget{sga3:xv:theorem:8-15}Theorem 8.15}{Let \(S\) be a locally noetherian
prescheme and let \(G\) be an \(S\)-group prescheme smooth and of
finite type over \(S\). The following conditions are equivalent:
\begin{enumerate}
\renewcommand{\labelenumi}{\textup{\roman{enumi})}}
\item The \(S\)-functor \(\mathcal T_M\), whose \(S'\)-points are
the maximal tori of \(G_{S'}\) (Exposé \SGARefLink{sga3:xii:1.3}{XII~1.3}), is representable.
\item The functor \(\mathcal T_M\) is represented by an
\(S\)-prescheme smooth and quasi-projective over \(S\), with affine
fibres.
\item The group \(G\) admits a maximal torus locally for the étale
topology.
\item The group \(G\) admits a maximal torus locally for the
faithfully flat topology.
\item\SGAInlineRefTarget{sga3:xv:theorem:8-15:item:v} The group \(G\) has property \(\mathrm{AT}\), and the
reductive rank of its fibres is a locally constant function on
\(S\).
\end{enumerate}}
\SGARefTarget{sga3:xv:theorem:8:15}

\SGAProof{Proof}
\textup{(ii)}\(\Rightarrow\)\textup{(i)} is clear.

\textup{(i)}\(\Rightarrow\)\textup{(iii)}.
Since \(G\) is of finite presentation over \(S\), \SGARefLink{sga3:xv:corollary:6-3}{Corollary 6.3}
shows that \(\mathcal T_M\) commutes with filtered inductive limits
of rings. If representable, it is therefore locally of finite
presentation (EGA IV~8.14). Moreover, \(\mathcal T_M\) is formally
smooth (Exposé XI~2.1 bis). Hence its representing prescheme is
smooth over \(S\), and \textup{(iii)} follows from Hensel's lemma
(Exposé \SGARefLink{sga3:xi:proposition:1:10}{XI~1.10}).

\textup{(iii)}\(\Rightarrow\)\textup{(iv)} is clear.

\textup{(iv)}\(\Rightarrow\)\textup{(v)}.
Let \(s\in S\). There is an \(S\)-prescheme \(S'\), flat over \(S\)
and with image containing \(s\), such that \(G_{S'}\) has a maximal
torus \(T'\). If \(s'\in S'\) lies over \(s\), the reductive rank of
the fibres of \(G_{S'}\) is constant on
\(\operatorname{Spec}\mathcal O_{S',s'}\). Thus the reductive rank
of the fibres of \(G\) is constant on its image
\(\operatorname{Spec}\mathcal O_{S,s}\)
(EGA IV~2.3.4(ii)). Since this rank is a locally constructible
function on \(S\) (6.3 bis), it is constant on a neighbourhood of
\(s\) (EGA IV~1.10.1).

To prove property \(\mathrm{AT}\), let \(S_1\) be an \(S\)-scheme
which is the spectrum of a discrete valuation ring, whose closed
point \(s_1\) maps to \(s\). The prescheme
\[
  S'_1=S_1\times_S S'
\]
is faithfully flat over \(S_1\), and \(G_{S'_1}\) has a maximal
torus. Let \(A\) and \(A'\) be the rings of \(S_1\) and \(S'_1\),
respectively. Writing \(A'\) as the filtered inductive limit of its
finitely generated \(A\)-subalgebras and applying \SGARefLink{sga3:xv:corollary:6-3}{Corollary 6.3},
we obtain an \(S_1\)-scheme \(S''_1\) such that
\(G_{S''_1}\) has a maximal torus and
\(S''_1\to S_1\) is of finite type and surjective. By
EGA II~7.1.9, we may reduce to the case in which \(S'_1\) is the
spectrum of a discrete valuation ring. Then \(G_{S'_1}\), and
hence \(G_{S_1}\), has property \(\mathrm{AT}\). Since this holds
for every such \(S_1\), \(G\) has property \(\mathrm{AT}\).

\textup{(v)}\(\Rightarrow\)\textup{(i)}.
By \SGARefLink{sga3:xv:theorem:8:8}{Theorem 8.9} the functor \(\mathcal T\) of sub-tori of \(G\) is
representable. If the reductive rank of \(G\) is the constant
function \(r\), then \(\mathcal T_M\) is represented by the
open-and-closed subprescheme of \(\mathcal T\) representing the
subfunctor of tori of rank \(r\).

\textup{(iii)}\(\Rightarrow\)\textup{(ii)}.
Under condition \textup{(iii)}, the results of Exposé \SGARefLink{sga3:xii:7.1}{XII~7.1}
identify \(\mathcal T_M\) canonically with the functor of Cartan
subgroups of \(G\). The assertion follows from Theorem 7.3(i).

\medskip\noindent\textbf{Remark 8.16.}
The prescheme \(\mathcal T_M\) of maximal tori of \(G\) can in fact
be shown to be affine over \(S\),\footnote{Cf. M. Raynaud,
\emph{Faisceaux amples sur les schémas en groupes et les espaces
homogènes} (thesis, forthcoming).} generalizing Exposé \SGARefLink{sga3:xii:5.4}{XII~5.4}.
\SGARefTarget{sga3:xv:remark:8:16}

\SGAStatement{Corollary 8.17}{Let \(S\) be a prescheme and let
\(G\) be an \(S\)-group prescheme smooth and of finite presentation
over \(S\). If the abelian and reductive ranks of the fibres of
\(G\) are locally constant functions on \(S\), then \(G\) satisfies
the equivalent conditions \textup{(i)}--\textup{(iv)} of
\SGARefLink{sga3:xv:theorem:8-15}{Theorem 8.15}.}
\SGARefTarget{sga3:xv:corollary:8:17}

We may suppose that the abelian and reductive ranks of the fibres
of \(G\) are constant. Proceeding as in \SGARefLink{sga3:xv:corollary:8:11}{Corollary 8.11} and using
\SGARefLink{sga3:xv:corollary:6:3:bis}{6.3 bis}, we reduce to the case in which \(S\) is noetherian. Then
\(G\) has property \(\mathrm{AT}\) by \SGARefLink{sga3:xv:proposition:8:7}{Proposition 8.7}, and the
conclusion follows from \SGARefLink{sga3:xv:theorem:8-15:item:v}{Theorem 8.15(v)}.

\SGAStatement{\SGAInlineRefTarget{sga3:xv:corollary:8-18}Corollary 8.18}{Let \(S\) be a prescheme and let
\(G\) be an \(S\)-group prescheme smooth over \(S\). The following
conditions are equivalent:
\begin{enumerate}
\renewcommand{\labelenumi}{\textup{\alph{enumi})}}
\item The unipotent rank \(\rho_u\) and the abelian rank
\(\rho_{\mathrm{ab}}\) of the fibres of \(G\) are locally constant
functions on \(S\).
\item The unipotent rank \(\rho_u\) is locally constant, and \(G\)
admits maximal tori locally for the fpqc topology.
\item The reductive rank \(\rho_r\) and the abelian rank
\(\rho_{\mathrm{ab}}\) of the fibres of \(G\) are locally constant
functions on \(S\).
\end{enumerate}}
\SGARefTarget{sga3:xv:corollary:8:18}

\medskip\noindent\textbf{Remark 8.19.}
Under the hypotheses of \SGARefLink{sga3:xv:corollary:8-18}{Corollary 8.18}, a more precise argument
using lower semicontinuity of the abelian rank, announced in
Exposé \SGARefLink{sga3:x:remark:8-7}{X~8.7}, shows that if two of the three ranks
\(\rho_u,\rho_r,\rho_{\mathrm{ab}}\) are locally constant, then so
is the third.
\SGARefTarget{sga3:xv:remark:8:19}

\SGAProof{Proof of \SGARefLink{sga3:xv:corollary:8-18}{Corollary 8.18}}
After replacing \(G\), if necessary, by its identity component, we
may suppose that \(G\) is of finite presentation over \(S\)
(Exposé VI\textsubscript{B}~\SGARefLink{par:sga3-VIA-5.3.3}{5.3.3}).

\textup{(a)}\(\Rightarrow\)\textup{(c)}.
Let \(s\in S\). Since \(\rho_{\mathrm{ab}}\) is locally constant,
\SGARefLink{sga3:xv:corollary:8:11}{Corollary 8.11} allows us, after an étale extension covering \(S\),
to suppose that \(G\) has a sub-torus \(T\) whose fibre \(T_s\) is
a maximal torus of \(G_s\). Put
\[
  C=\operatorname{Centr}_G(T).
\]
This is a subgroup prescheme of \(G\), smooth over \(S\), with
connected fibres. For every \(t\in S\), the group \(C_t\) contains
a Cartan subgroup \(C'_t\) of \(G_t\). After shrinking \(S\), we
may suppose that \(\rho_u\), \(\rho_{\mathrm{ab}}\), and the
relative dimension of \(C\) over \(S\) are constant. We then have
\[
\begin{aligned}
  \dim C_s=\dim C_t
  &\geq \dim C'_t\\
  &\geq \rho_u(t)+\rho_{\mathrm{ab}}(t)+\dim T_t\\
  &=\rho_u(s)+\rho_{\mathrm{ab}}(s)+\dim T_s\\
  &=\rho_\nu(s)=\dim C_s.
\end{aligned}
\]
It follows that \(C_t=C'_t\). Thus \(T\) is a maximal torus of
\(G\), and in particular \(\rho_r(t)=\rho_r(s)\).

\textup{(c)}\(\Rightarrow\)\textup{(b)} follows from
\SGARefLink{sga3:xv:corollary:8:17}{Corollary 8.17}.

\textup{(b)}\(\Rightarrow\)\textup{(a)}.
Since \(G\) admits maximal tori locally for the fpqc topology, it
also admits Cartan subgroups locally for that topology. Hence, by
Exposé \SGARefLink{sga3:xii:7.3}{XII~7.3}, the nilpotent rank
\[
  \rho_\nu=\rho_u+\rho_r+\rho_{\mathrm{ab}}
\]
is locally constant. The existence of maximal tori makes
\(\rho_r\) locally constant, while \(\rho_u\) is locally constant
by hypothesis; therefore \(\rho_{\mathrm{ab}}\) is locally
constant.\footnote{The final sentence of the printed French
interchanges \(\rho_u\) and \(\rho_{\mathrm{ab}}\). The displayed
identity and condition \textup{(b)} give the implication in the
form written here.}
